\section{Guarantees and Correctness Properties}
\label{sec:properties}

The preceding sections specify the Bailout Protocol as a deterministic state machine, an escalation mechanism operating across a recursive agent tree, and an implementation realized within the \bxthree{} three-layer architecture. This section translates that specification into three formal correctness properties — \textbf{Safety} (no autonomous collapse), \textbf{Liveness} (eventual human contact), and \textbf{Accountability} (human always identifiable) — and provides proof sketches for each under the stated structural assumptions. Together, these properties constitute the operational meaning of the upstream accountability guarantee for the Bailout Protocol.

\bigskip
% --------------------------------------------------------------
\subsection{Safety: No Autonomous Collapse}
\label{sec:safety}

\textbf{Statement (Safety).} For any \bxthree{} node $n$, no autonomous execution path exists that remains in a non-\texttt{IDLE} state indefinitely without reaching the human anchor $h(n)$.

\bigskip
\noindent
Formally: $\forall n \in \text{nodes},\; \forall s \in Q \setminus \{\text{IDLE}\},\; \text{state}(n) = s \implies \exists k \geq 0,\; \text{step}^k(n, s) = h(n)$.

\bigskip
\noindent
\textbf{Structural assumptions.}
\begin{itemize}\itemsep=0pt
\item[\namedlabel{asm:layer}]{(A1)} Every node $n$ carries a \purpose{} layer, a \bounds{} layer, and a \fact{} layer with the enforcement properties described in Section~\ref{sec:implementation}.
\item[\namedlabel{asm:immutability}]{(A2)} The human anchor $h(n)$ and the parent pointer $\alpha(n)$ are immutable for the operational lifetime of $n$.
\item[\namedlabel{asm:fact-exclusive}]{(A3)} The \fact{} layer holds exclusive write authority over the authorization ledger; no machine actor may commit a ledger entry without \fact{}-layer approval.
\item[\namedlabel{asm:transition-complete}]{(A4)} The state machine transition relation $\rightarrow$ is complete for all $(q, \sigma) \in Q \times \Sigma$ with respect to the escalation path: every non-\texttt{IDLE} state has at least one defined outgoing transition.
\item[\namedlabel{asm:timeout}]{(A5)} All timeout parameters ($T_{\text{bounds}}$, $T_{\text{propagate}}$, $T_{\text{human}}$) are positive finite values.
\end{itemize}

\bigskip
\noindent
\textbf{Proof sketch.} The proof proceeds by structural induction over the state space $Q$ of the Bailout Protocol state machine.

\medskip
\noindent
\textit{Base case.} The state \texttt{BOUNDED} is a non-\texttt{IDLE} state. By the state machine specification (T1–T3), \texttt{BOUNDED} has two defined exit transitions: T2 returns to \texttt{IDLE} when a \fact-layer-approved resolution is produced within $T_{\text{bounds}}$, and T3 advances to \texttt{BAIL\_PENDING} when $T_{\text{bounds}}$ expires or all resolution paths are exhausted. In the first case, the protocol resolves without escalation. In the second case, the state machine proceeds to \texttt{BAIL\_PENDING}, which by INV-2 (Section~\ref{sec:state-machine-fact-layer}) immediately and atomically transitions to \texttt{ESCALATING} via T4. There is no path from \texttt{BOUNDED} to a stall state; the only resolution is either local resolution (returning to \texttt{IDLE}) or escalation to $h(n)$.

\medskip
\noindent
\textit{Inductive step.} Assume for some state $s \in Q \setminus \{\text{IDLE}, \text{BOUNDED}\}$ that every execution path from $s$ either returns to \texttt{IDLE} or reaches $h(n)$. Consider state \texttt{ESCALATING}. By the Escalate algorithm (Algorithm~1, Section~\ref{sec:bailout}), the node transmits the exception context to its parent via all active functional pathways and awaits acknowledgment. If any pathway returns \texttt{ACK} within $T_{\text{propagate}}$, the state machine transitions to \texttt{HUMAN\_ACKNOWLEDGED} (T5) and then to \texttt{RESOLVED} (T6), closing the protocol at $h(n)$ with a full Forensic Ledger entry. If no pathway acknowledges within $T_{\text{propagate}}$, the protocol times out (T7) and enters \texttt{CIRCUIT\_BROKEN}, which by the timeout handling specification (Section~\ref{sec:bailout}) transitions directly to a state in which $h(n)$ is contacted via an out-of-band pathway (direct communication, backup channel, or institutional escalation). In all cases, the execution path from \texttt{ESCALATING} terminates at $h(n)$ — no stall state is reachable from \texttt{ESCALATING}.

\medskip
\noindent
\textit{No autonomous collapse.} An autonomous collapse would require a non-\texttt{IDLE} execution path that indefinitely avoids both local resolution (returning to \texttt{IDLE}) and escalation to $h(n)$. By the inductive argument above, every non-\texttt{IDLE} state either returns to \texttt{IDLE} via an approved resolution, transitions to a state that has $h(n)$ as its defined terminus, or times out and triggers an out-of-band human contact procedure. No state in $Q$ has an undefined exit transition that would produce a stall. Therefore no autonomous execution path can remain in a non-\texttt{IDLE} state indefinitely without reaching $h(n)$. $\square$

\medskip
The Safety property establishes that the Bailout Protocol is structurally incapable of producing the Intercept failure mode described in Section~\ref{sec:properties}: a machine actor cannot absorb an exception and continue autonomous operation because every non-\texttt{IDLE} state has a defined path to $h(n)$, and no machine actor has the authority to close the protocol without a \purpose-layer-authorized directive.

\bigskip
% --------------------------------------------------------------
\subsection{Liveness: Eventual Human Contact}
\label{sec:liveness}

\textbf{Statement (Liveness).} For any \bxthree{} node $n$ that enters \texttt{ESCALATING} state, $h(n)$ will receive the exception context within a bounded time from the moment of entry, provided the physical communication infrastructure to $h(n)$ is operational.

\bigskip
\noindent
Formally: $\forall n \in \text{nodes},\; \text{state}(n) = \text{ESCALATING} \implies \exists t \leq T_{\text{liveness}},\; \text{received}(h(n), \text{context}(n))$, where $T_{\text{liveness}} = T_{\text{propagate}} + T_{\text{timeout}} + T_{\text{oob}}$ and $T_{\text{oob}}$ is the maximum out-of-band contact latency.

\bigskip
\noindent
\textbf{Structural assumptions (in addition to A1–A5).}
\begin{itemize}\itemsep=0pt
\item[\namedlabel{asm:pathway}]{(A6)} At least one functional pathway to $h(n)$ is operational at the time of entry to \texttt{ESCALATING}, or the physical communication infrastructure enabling out-of-band contact is operational.
\item[\namedlabel{asm:human-available}]{(A7)} $h(n)$ is reachable (not permanently offline) within the time bounds defined by the protocol's timeout parameters.
\end{itemize}

\bigskip
\noindent
\textbf{Proof sketch.} The proof proceeds by case analysis on the communication pathway status at the time the node enters \texttt{ESCALATING}.

\medskip
\noindent
\textit{Case 1: At least one functional pathway is operational.} The node invokes the Escalate algorithm, which transmits the exception context simultaneously across all active functional pathways (Algorithm~1, lines 5–11). For each pathway $p$, the algorithm awaits \texttt{ACK}$(p, n)$ for a duration of $T_{\text{propagate}}$. If any pathway returns \texttt{ACK} within this window, the exception context is received by $h(n)$ at time $t_{\text{receive}} \leq T_{\text{propagate}} < T_{\text{liveness}}$. This satisfies the liveness condition.

\medskip
\noindent
\textit{Case 2: No functional pathway is operational ($T_{\text{propagate}}$ expires).} The Bailout Monitor records the timeout event and the state machine transitions from \texttt{ESCALATING} to \texttt{CIRCUIT\_BROKEN} via T7. The circuit-break handling procedure activates out-of-band contact with $h(n)$ — a direct communication mechanism (human telephone, institutional on-call, or equivalent) that operates outside the functional pathway infrastructure. The protocol specifies that out-of-band contact completes within $T_{\text{oob}}$, a provisioned finite parameter. The exception context is therefore received by $h(n)$ at time $t_{\text{receive}} \leq T_{\text{propagate}} + T_{\text{oob}} \leq T_{\text{liveness}}$.

\medskip
\noindent
\textit{Case 3: $h(n)$ does not acknowledge within $T_{\text{human}}$ after receiving the context.} This case does not violate the liveness condition as stated, because liveness is defined as receipt of the exception context by $h(n)$, not as receipt of a resolution directive by $n$. Once $h(n)$ has received the context, the protocol has satisfied its liveness obligation. The subsequent timeout at the protocol level (the $T_{\text{human}}$ expiration) triggers a \texttt{CIRCUIT\_BROKEN} state that logs the non-response and maintains the node in a quiescent state until $h(n)$ responds, but this does not constitute a liveness failure of the protocol itself.

\medskip
In all cases, the exception context reaches $h(n)$ within $T_{\text{liveness}}$, a bound determined entirely by the timeout parameters provisioned at node initialization and the out-of-band contact latency. $\square$

\medskip
The Liveness property does not require that $h(n)$ responds within a bounded time — that would impose a requirement on a human actor that is outside the system's control. It requires only that $h(n)$ receives the exception context within a bounded time, which is a property the protocol can guarantee structurally through its multi-pathway escalation and out-of-band fallback mechanisms.

\bigskip
% --------------------------------------------------------------
\subsection{Accountability: Human Always Identifiable}
\label{sec:accountability}

\textbf{Statement (Accountability).} For any exception state processed by the Bailout Protocol at any node $n$, the human anchor $h(n)$ is uniquely and immutably identifiable in the Forensic Ledger, and no protocol run can terminate without a \purpose-layer-authorized directive attributable to $h(n)$.

\bigskip
\noindent
Formally: $\forall n \in \text{nodes},\; \forall \text{protocol run } r \text{ at } n,\; \text{received}(h(n), \text{context}(n)) \implies \exists!\, h(n) \in \text{ForensicLedger}(r)$ and $\text{terminated}(r) \implies \exists \text{ledger entry } \ell \in \text{ForensicLedger}(r)$ such that $\ell.\text{directive} \in \{\texttt{ACK}, \texttt{RESOLVE}, \texttt{REJECT}\}$ and $\ell.\text{authorizer} = h(n)$.

\bigskip
\noindent
\textbf{Structural assumptions (in addition to A1–A7).}
\begin{itemize}\itemsep=0pt
\item[\namedlabel{asm:immutable-anchor}]{(A8)} $h(n)$ is established at node provisioning and immutably recorded in the authorization ledger; no machine actor may modify this assignment.
\item[\namedlabel{asm:directive-exclusivity}]{(A9)} Only the \purpose{} layer of $h(n)$ may authorize a termination directive; no machine actor may issue, simulate, or substitute a directive on behalf of $h(n)$.
\item[\namedlabel{asm:ledger-immutability}]{(A10)} The Forensic Ledger is append-only with SHA-256 hash-chain integrity, preventing retrospective modification of any entry after commit.
\end{itemize}

\bigskip
\noindent
\textbf{Proof sketch.} The proof has two components: (i) uniqueness and immutability of $h(n)$ identification, and (ii) directive non-repudiation for every protocol termination.

\medskip
\noindent
\textit{Component 1: Unique and immutable identification of $h(n)$.} At node provisioning (Section~\ref{sec:purpose-layer}, Stage (i)), the \purpose{} layer registers $h(n)$ with the \fact{} Layer's Pathway Registry and records the anchor assignment in the authorization ledger. This entry is committed by the \fact{} layer with exclusive write authority (A3). By (A8), $h(n)$ is immutably stored in this entry and cannot be modified by any machine actor — including the node's own Bounds Engine, any parent node, or any recursively spawned subordinate. When the Escalate algorithm is invoked, the algorithm retrieves $h(n)$ directly from the authorization ledger (Algorithm~1, line 3), not from any runtime-computed or machine-maintained reference. The retrieval is therefore traceable to the original provisioning entry, and any subsequent protocol run at node $n$ preserves this identification in its Forensic Ledger entries. No two nodes share the same $h(n)$ unless explicitly co-provisioned with shared accountability; in any case, the identification in the ledger is unique to the provisioning record.

\medskip
\noindent
\textit{Component 2: Directive non-repudiation for every termination.} A protocol run terminates only through one of the following transitions: T2 (\texttt{IDLE} via local resolution), T5 (\texttt{HUMAN\_ACKNOWLEDGED}), T6 (\texttt{RESOLVED}), or the timeout path T7 leading to \texttt{CIRCUIT\_BROKEN}. 

\begin{itemize}\itemsep=0pt
\item[T2:] The protocol terminates at \texttt{IDLE} via a \fact-layer-approved resolution. The \fact{} layer's authorization ledger gate (Section~\ref{sec:fact-layer}, enforcement mechanism ii) requires that the resolution be recorded as a ledger entry bearing \purpose{}-layer authorization. Because only the \purpose{} layer can authorize a ledger entry (A9), and the \purpose{} layer belongs to $h(n)$ by construction, the termination directive is attributable to $h(n)$ with non-repudiation.

\item[T5/T6:] The protocol terminates at \texttt{HUMAN\_ACKNOWLEDGED} or \texttt{RESOLVED} upon receipt of a directive from $h(n)$ ($\texttt{ACK}$, $\texttt{RESOLVE}$, or $\texttt{REJECT}$). The directive is recorded in the Forensic Ledger as a \purpose-layer-authorized entry (Section~\ref{sec:purpose-layer}, Stage (iii)), with the authorizer field set to $h(n)$. The SHA-256 hash chain (A10) prevents retrospective modification of this entry.

\item[T7:] The protocol enters \texttt{CIRCUIT\_BROKEN} after $T_{\text{propagate}}$ expires and all functional pathways have failed. The circuit-break handling activates out-of-band contact with $h(n)$ and logs the event. The protocol run does not close until $h(n)$ issues a directive; the node remains in \texttt{CIRCUIT\_BROKEN} until then. The eventual termination therefore still requires a \purpose-layer-authorized directive attributable to $h(n)$, and the \texttt{CIRCUIT\_BROKEN} state itself is recorded as a Forensic Ledger entry with full diagnostic context.
\end{itemize}

In all termination paths, the protocol run closes with a ledger entry bearing a \purpose-layer-authorized directive from $h(n)$. No termination path exists that closes without this entry, because the \fact{} layer's authorization ledger gate (A3) blocks any attempt to close a protocol run without one. $\square$

\medskip
The Accountability property establishes that the Bailout Protocol satisfies the upstream accountability guarantee at the level of individual protocol runs: every exception that the protocol processes is traceable to a specific human principal, and no protocol run can terminate in a way that obscures or denies the human's involvement. The combination of immutable $h(n)$ registration, \purpose{}-layer-exclusive directive authority, and append-only hash-chained Forensic Ledger entries makes both identification and non-repudiation structurally guaranteed rather than relying on behavioral assumptions about machine actors.

\bigskip
% --------------------------------------------------------------
\subsection{Collective Correctness of the Five Pillars}
\label{sec:collective-correctness}

The Safety, Liveness, and Accountability properties established above describe the behavior of the Bailout Protocol in isolation. However, these properties operate within the context of the full \bxthree{} framework, where the five pillars function as an interdependent accountability system.

The Bailout Protocol's correctness properties rely critically on the structural assumptions listed in (A1)–(A10) above. These assumptions are themselves guaranteed by the other four pillars of \bxthree{}:

\begin{itemize}\itemsep=0pt
\item \textbf{Loop Isolation} ensures that no node can corrupt another's execution state, which is necessary for the state machine invariant INV-1 to hold (Section~\ref{sec:state-machine-fact-layer}).

\item \textbf{Recursive Spawning with Immutable Parent Pointers} establishes the parent chain structure that the Escalate algorithm traverses, and guarantees assumption (A2) — the immutability of $\alpha(n)$ and $h(n)$.

\item \textbf{Spatial Firewall} prevents cross-domain state contamination that could compromise the \fact{} Layer's enforcement of the transition relation, protecting the atomicity of state transitions.

\item \textbf{Sandbox Gate} ensures that the operating range $R(n)$ provisioned for each node is enforced at its boundary, preventing unauthorized state transitions that could bypass the trigger condition.
\end{itemize}

The five pillars collectively establish the preconditions under which the Bailout Protocol's correctness properties hold. A system that implements the Bailout Protocol without the other four pillars cannot guarantee assumptions (A1)–(A10), and therefore cannot guarantee Safety, Liveness, and Accountability as stated. The upstream accountability guarantee — that \bxthree{} systems fail upward into human consciousness, never downward into autonomous chaos — is thus a property of the complete five-pillar system, not of the Bailout Protocol alone.